\subsection{Theoretical Foundations of Cross-Modal Forensic Alignment}
\label{sec:cmfa_theory}

We formalize the three primary detection signals leveraged by the Cross-Modal Forensic Alignment (CMFA) module. The efficacy of CMFA lies in its cross-signal redundancy, capturing discrepancies across color space, structural font metrics, and high-frequency compression statistics.

\subsubsection{Limitations of Traditional Error Level Analysis (ELA) on Mobile Interfaces}
Traditional ELA is widely utilized to identify spliced regions in compressed raster images. The underlying mechanism relies on the mathematical properties of the $8 \times 8$ Discrete Cosine Transform (DCT) grid boundaries established during initial JPEG compression. When a JPEG image is modified and recompressed, the mismatch between the original grid and the newly aligned grid results in non-uniform error distributions.

However, mobile user interfaces (UIs) present a unique forensic challenge. Modern mobile operating systems capture screenshots natively in lossless Portable Network Graphics (PNG) format. Consequently, the input screenshot $I$ lacks any prior lossy compression artifacts:
\begin{equation}
    I = I_{\text{lossless}} \implies E_{\text{original}} = 0
\end{equation}
If traditional ELA is applied directly by applying a secondary JPEG compression layer at quality $Q=90$, the resulting error map $E = |I - I_c|$ yields a uniform error distribution across the entire image because the DCT coefficients are quantized uniformly for the first time. The error is globally distributed and lacks high-contrast local variance, rendering traditional ELA uninformative.

To overcome this, Lumint introduces \textit{grid-density ELA}. Instead of assessing raw pixel difference directly, we partition the error map $E$ into a $12 \times 12$ grid of blocks $\{G_{ij}\}$. We define a local error density operator. Let $\mathcal{H}(G_{ij})$ be an indicator function mapping a block to a hotspot if its mean error exceeds a threshold $\tau_{\text{ela}} = 127$:
\begin{equation}
    \mathcal{H}(G_{ij}) = \mathbb{I}(\text{mean}(G_{ij}) > \tau_{\text{ela}})
\end{equation}
A \textit{localized hotspot} is defined as a block $G_{ij}$ where $\mathcal{H}(G_{ij}) = 1$ surrounded by a neighborhood of low error $\mathcal{N}(G_{ij})$ such that $\forall G_{kl} \in \mathcal{N}(G_{ij})$, $\text{mean}(G_{kl}) < \tau_{\text{ela}}/2$. Conversely, a \textit{global compression artifact} exhibits uniform error:
\begin{equation}
    \forall i,j \quad \text{mean}(G_{ij}) \approx \mu_{\text{global}} \pm \epsilon
\end{equation}
By calculating the density $f_3 = \sum_{i,j} \mathcal{H}(G_{ij}) / 144$, we isolate localized modifications (such as splicing transaction amounts) from uniform background noise.

\subsubsection{Font Splicing Statistical Model}
Digital payment interfaces render text elements using specific system font rendering engines (e.g., San Francisco on iOS, Roboto/Inter on Android) with rigid layout configurations. Let $H = \{h_1, h_2, \dots, h_N\}$ represent the set of detected bounding box heights for alphanumeric characters within the screenshot.

In a genuine screenshot, all text elements are generated by a single app-specific font renderer $\mathcal{R}_{\text{app}}$. Under this homogeneous rendering process, character heights conform to a tightly bound normal distribution:
\begin{equation}
    h_i \sim \mathcal{N}(\mu_{\text{app}}, \sigma^2_{\text{renderer}})
\end{equation}
where the variance $\sigma^2_{\text{renderer}}$ is strictly bounded by the typographical constraints of the app template:
\begin{equation}
    \sigma^2_{\text{renderer}} < 5.0
\end{equation}
In a forged screenshot generated via template overlay (where a fraudster splices a text element from a web browser or image editor), text contours originate from multiple distinct rendering engines $\mathcal{R}_1, \dots, \mathcal{R}_k$, each governed by differing mean heights $\mu_1, \dots, \mu_k$. The combined height distribution of the forged screenshot becomes a mixture model:
\begin{equation}
    P(h) = \sum_{k=1}^K \pi_k \mathcal{N}(\mu_k, \sigma^2_k)
\end{equation}
The total variance of this mixture distribution is mathematically inflated:
\begin{equation}
    \sigma^2_{\text{forged}} = \sum_{k=1}^K \pi_k \sigma^2_k + \sum_{k=1}^K \pi_k (\mu_k - \bar{\mu})^2 = \sigma^2_{\text{genuine}} + \text{Var}(\{\mu_k\})
\end{equation}
CMFA exploits this variance inflation. Empirical analysis of the UPI-FraudBench dataset demonstrates that the 99th percentile of genuine sample height variance is $108.0$. We set our threshold at $\tau_{\text{var}} = 110.0$; any height variance $\sigma^2 > 110.0$ indicates multi-source font splicing.

\subsubsection{Brand Palette Attack Resilience Analysis}
Consider an adversary $\mathcal{A}$ attempting to deceive the color authenticity detector. Let $T_a$ be the target brand color vector. The adversary's goal is to modify the screenshot such that the extracted dominant colors $D$ align with the brand template $T_a$, while adhering to a maximum pixel perturbation budget $\epsilon_{\text{color}}$ in the RGB color space:
\begin{equation}
    \|I_{\text{forged}}(x,y) - I_{\text{genuine}}(x,y)\|_2 \le \epsilon_{\text{color}}
\end{equation}
To minimize the brand palette distance $\delta_{\text{color}} = \min_{\mathbf{d} \in D} \|\mathbf{d} - T_a\|_2$ below the detection threshold $\tau_{\text{color}} = 85.0$, the adversary must shift the dominant colors of the spliced region. Since dominant colors are extracted via histogram bucketing with a bin width of $32$, any perturbation smaller than $32$ units per channel fails to shift the pixel into the correct dominant bin. Therefore, the adversary requires:
\begin{equation}
    \epsilon_{\text{color}} \ge 32 \sqrt{3} \approx 55.4
\end{equation}
This shift introduces high-frequency edges along the spliced boundaries. Specifically, a pixel value adjustment of $\ge 32$ units per channel creates a localized boundary discontinuity. When the forged screenshot is compressed at $Q=90$ during ELA computation, the boundary discontinuity acts as a high-frequency step function. The JPEG quantization step amplifies the error at this boundary, causing the local ELA density $f_3$ to exceed the threshold $\tau_{\text{ela}}$, triggering detection. This mathematical coupling between color perturbation and compression noise proves that no adversary can simultaneously bypass the color distance and ELA detectors.
